\appendix

\section{Per-Experiment Parameters}
\label{app:exp_params}

\tabref{tab:exp_params} collects the per-experiment values for reproducibility. Values shared across experiments---propagation horizon \SI{24}{\hour}, $K=240$ time steps, minimum elevation $\alpha_{\min} = \SI{10}{\degree}$, AdamW with learning rate $\eta=10^{-2}$---are omitted. The ground grid is $36 \times 72$ except where a regional target set is specified.

\begin{table}[!htb]
\centering
\caption{Per-experiment parameters for the experiments in \secref{sec:experiments}.}
\label{tab:exp_params}
\begin{tabular}{lccccccc}
\toprule
\textbf{Experiment} & $N$ & \textbf{DoF} & $J$ & $\tau_{\text{cov}}$ [\si{\degree}] & $\tau_{\text{rev}}$ [\si{\degree}] & $\beta$ [\si{\minute}] & $\lambda$ \\
\midrule
Exp 1 (toy)           &  2 &  2 & 2592 & 2.0 & 2.0 & 10 & 2.0 \\
Exp 2 (Walker)        & 24 & 30 & 2592 & 2.0 & 2.0 & 10 & 0.1 \\
Exp 3 (Europe)        &  4 & 12 &  500 & 2.0 & 2.0 & 10 & 1.0 \\
Baselines (SA/GA/DE)  & 24 & 30 & 2592 & --- & --- & ---& 0.1 \\
Hyperparam tuning     & 24 & 30 & 2592 & 2.0 & 2.0 & 10 & 0.1 \\
\bottomrule
\end{tabular}
\end{table}

\section{Hyperparameter Tuning}
\label{app:hyperparam_grid}

This appendix provides the complete experimental setup and raw results for hyperparameter tuning described in \secref{sec:tightness}.

\paragraph{Constellation configuration.}
The tuning uses a 24-satellite constellation ($6 \times 4$) at \SI{550}{\kilo\metre} altitude with eccentricity $\approx 0$, inclination fixed at \SI{60}{\degree}, and RAAN and mean anomaly as free parameters.
RAAN is shared within each plane, giving 30 effective degrees of freedom (6 unique RAANs + 24 mean anomalies).

\paragraph{Optimization settings.}
Each of the four initial configurations is optimized for 800 iterations with AdamW ($\eta = 10^{-2}$), propagation over \SI{24}{\hour} discretized into $K = 240$ time steps, a ground grid of $36 \times 72 = 2592$ points spanning $\pm$\SI{70}{\degree} latitude with $\cos(\text{lat})$ area weights, minimum elevation $\alpha_{\min} = \SI{10}{\degree}$, sigmoid temperature $\tau = \SI{2}{\degree}$, LogSumExp temperature $\beta = \SI{10}{\minute}$, revisit weight $\lambda = 0.1$, mean revisit reduction, and GMST randomization enabled.
The optimization uses the \emph{default} relaxation parameters, not the calibrated ones---the goal is to obtain four diverse final solutions whose relative quality is known from the hard metrics.

\paragraph{Initial RAAN configurations.}
The four configurations are: near-uniform $[0, 60, 120, 180, 240, 300]$\SI{}{\degree}, moderate $[0, 30, 120, 200, 210, 300]$\SI{}{\degree}, clustered $[0, 10, 20, 30, 40, 50]$\SI{}{\degree}, and two-cluster $[0, 5, 180, 185, 270, 275]$\SI{}{\degree}.
All use the same random mean anomaly initialization (NumPy \texttt{RandomState(42)}, uniform on $[0, 360)$\SI{}{\degree}).

\paragraph{Grid search parameters.}
The grid search evaluates the relaxed loss of all four final solutions under each combination of: coverage softness $\tau_{\text{cov}} \in \{1.0, 1.5, 2.0, 2.5, 3.0, 4.0, 5.0\}$\SI{}{\degree}, revisit softness $\tau_{\text{rev}} \in \{0.1, 0.25, 0.5, 0.75, 1.0, 1.5, 2.0\}$\SI{}{\degree}, LogSumExp temperature $\beta \in \{3.0, 5.0, 7.5, 10.0, 15.0, 20.0\}$\SI{}{\minute}, and revisit weight $\lambda \in \{0.05, 0.1, 0.2, 0.5, 1.0, 2.0\}$.
This gives $7 \times 7 \times 6 \times 6 = 1764$ combinations, each requiring 4 gradient-free forward passes (one per solution), for 7056 evaluations total.

\paragraph{Validity criterion.}
Each of the four optimizations converges to a different local minimum. A hyperparameter combination is \emph{valid} if the relaxed loss at these local minima preserves the same ordering as the hard loss---i.e., solutions that are better under the hard metrics should also be better under the relaxed loss. Since the hard metrics group the four solutions into two tiers (near-uniform and moderate are better; clustered and two-cluster are worse), validity requires:
\begin{equation}
    \max\bigl(\tilde{\mathcal{L}}_{\text{near-uniform}},\, \tilde{\mathcal{L}}_{\text{moderate}}\bigr) < \min\bigl(\tilde{\mathcal{L}}_{\text{clustered}},\, \tilde{\mathcal{L}}_{\text{two-cluster}}\bigr).
    \label{eq:validity}
\end{equation}
The \emph{margin} is the difference between the right- and left-hand sides.
Among valid combinations, we utilize the one that maximizes $\tau_{\text{cov}}$ (smoothest coverage gradients), then $\tau_{\text{rev}}$, with margin as tiebreaker.

\paragraph{Reproducibility.}
All experiments can be reproduced from the repository found at \url{https://github.com/shreeyam/differentiable_eo}.

\section{Equivalences among candidate relaxations}
\label{app:alternatives}

Several natural alternatives to R1, R2, and R4 reduce \emph{exactly} to the chosen relaxations up to reparameterisation or additive constants.
We record these identities because they explain why swapping in the alternative would not produce different optimizer dynamics.

\subsection{R1: tanh is a rescaled sigmoid}

For any $x \in \mathbb{R}$,
\begin{equation}
    \tfrac{1}{2}\bigl(1 + \tanh(x/2)\bigr)
    = \tfrac{1}{2}\left(1 + \frac{e^{x/2} - e^{-x/2}}{e^{x/2} + e^{-x/2}}\right)
    = \frac{e^{x/2}}{e^{x/2} + e^{-x/2}}
    = \frac{1}{1 + e^{-x}}
    = \sigma(x),
    \label{eq:tanh_sigmoid}
\end{equation}
so the rescaled tanh $\bigl(1 + \tanh((\alpha_{ij} - \alpha_{\min})/(2\tau))\bigr)/2$ is pointwise identical to the sigmoid in \eqref{eq:soft_visibility}.
Gradients, level sets, and every downstream quantity are unchanged.

\subsection{R2: Noisy-OR is the probabilistic t-conorm}

In fuzzy logic the \emph{probabilistic t-conorm}~\citep{klement_triangular_2000} of $n$ values $\{x_i\} \subset [0,1]$ is defined as $1 - \prod_i (1 - x_i)$, which is exactly the noisy-OR of \eqref{eq:noisy_or}.
Equivalently, if $\tilde{c}_{ij}(t_k)$ is interpreted as the independent detection probability of satellite $i$ at target $j$, then
\begin{equation}
    \Pr\!\left[\bigcup_i \text{sat } i \text{ covers } j\right]
    = 1 - \prod_i \Pr\!\left[\text{sat } i \text{ misses } j\right]
    = 1 - \prod_i (1 - \tilde{c}_{ij}(t_k)).
    \label{eq:noisy_or_derivation}
\end{equation}
No separate relaxation is needed: the noisy-OR operator is the unique member of the probabilistic t-conorm family that is smooth, bounded in $[0,1]$, and strictly increasing in each argument.

\subsection{R4: Mellowmax differs from LogSumExp by an additive constant}

The mellowmax operator was introduced by \citet{asadi_alternative_2017} as an alternative to the Boltzmann softmax used for value aggregation in $Q$-learning, where it restores a non-expansion property Boltzmann softmax lacks.
Despite the motivation, mellowmax is itself an instance of LogSumExp up to an additive constant:
\begin{equation}
    \operatorname{mm}_{\omega}(X)
    \;=\; \frac{1}{\omega}\log\!\left(\frac{1}{K}\sum_{k=1}^{K} e^{\omega x_k}\right).
    \label{eq:mellowmax}
\end{equation}
Setting $\omega = 1/\beta$ and expanding the $1/K$ factor gives
\begin{equation}
    \operatorname{mm}_{1/\beta}(X)
    = \beta \log\!\left(\sum_k e^{x_k/\beta}\right) - \beta \log K
    = \logsumexp_\beta(X) - \beta \log K,
    \label{eq:mellowmax_to_lse}
\end{equation}
so the two differ by the constant $\beta \log K$.
Since gradients are invariant under additive constants, the optimizer sees identical dynamics; using mellowmax here would be a relabeling of LogSumExp rather than a different relaxation.